\section{Conclusion}
We studied segment running time estimation with a known total duration, combining estimates shared across repeated intervals with current-trip adjustments. The base branch uses interval identities and context, the trip branch redistributes time within a relative bound, and final rescaling enforces the supplied total. Outputs remain nonnegative and sum-consistent, although the total constraint alone cannot guarantee interval accuracy.

Segment means remain competitive with very few labels. Ours has lower point-estimate MAE than direct prediction at 1\% and 5\% (by 0.579 and 0.075 seconds), but higher MAE at 10\%; paired runs favor Ours in eight, six, and zero of nine cases. Its long-interval MAE is higher at every budget. The validation-selected Ours+Direct ensemble has the lowest mean MAE at 5\% and 10\%, while segment means remain better at 1\%, but the ensemble is a separate predictor. These development results do not establish a uniform advantage for Ours. Equal label counts also yield different errors under different interval coverage. Independent totals, an unopened external evaluation, and data from other cities are needed to determine whether the allocation structure improves local estimates beyond these development comparisons.
